\section*{Chapter \ref{ch:m3:emissions-and-environmental-markets} --- Emissions and Environmental Markets}

\begin{solution}{exo:m3:emissions-and-environmental-markets:1}
It surrenders 1\,000 allowances, by 30 September of the year after the emissions; the 100 bought
cover the shortfall.
\end{solution}

\begin{solution}{exo:m3:emissions-and-environmental-markets:2}
At 1\,200 million it withholds 24\%, 288 million; at 900 million, the excess over 833 million, 67
million; at 350 million it releases 100 million.
\end{solution}

\begin{solution}{exo:m3:emissions-and-environmental-markets:3}
Under a tax the price stays and emissions fall with demand; under a cap the quantity stays at the cap
and the price falls, leaving a surplus that can depress prices for years unless a mechanism such as the
reserve reduces future supply.
\end{solution}

\begin{solution}{exo:m3:emissions-and-environmental-markets:4}
\qty{13.11}{EUR/t} with gas at \qty{20}{}; \qty{60.47}{EUR/t} with coal at 38\% and gas at
\qty{35}{}: the less efficient coal plant emits more per MWh, so a lower carbon price suffices.
\end{solution}

\begin{solution}{exo:m3:emissions-and-environmental-markets:5}
$0.202/0.55 = 0.367$~t per MWh; a 10-euro rise adds \qty{3.67}{EUR/MWh}.
\end{solution}

\begin{solution}{exo:m3:emissions-and-environmental-markets:6}
A \omterm{def:m3:emissions-and-environmental-markets:go}{guarantee of origin} transfers a claim about how a MWh was produced; the renewable MWh was produced
anyway and no cap binds the buyer. A surrendered allowance is cancelled within a cap: the total that
can be emitted is fixed, so each tonne covered is one less available to anyone else.
\end{solution}

\begin{solution}{exo:m3:emissions-and-environmental-markets:7}
Between 833 and 1\,096 million the ratio $(T - 833)/T$ rises with $T$, to 23.996\% at 1\,096
million; above it the intake is 24\% of $T$. At 1\,096 million the two pieces give 263.00 and 263.04
million: continuous within 0.04 million.
\end{solution}

\begin{solution}{exo:m3:emissions-and-environmental-markets:8}
A low price can mean the cap is being met cheaply (cheap abatement, low demand), not that it does not
bind: emissions cannot exceed the cap whatever the price. The price measures the marginal cost of
staying under the cap, and the surplus it leaves is banked for later years, when the cap is tighter.
\end{solution}

\begin{solution}{pb:m3:emissions-and-environmental-markets:1}
\textbf{1.} Coal $0.341/0.40 = 0.8525$~t per MWh; gas $0.202/0.55 = 0.3673$~t. \textbf{2.} $6 \times
0.8525 + 4 \times 0.3673 = 6.584$ million tonnes. \textbf{3.} About EUR~460.9 million. \textbf{4.} By
30 September of the year after. \textbf{5.} To fix the cost when it sells its power forward, and because
waiting leaves it exposed to the price for up to 21 months.
\textbf{6.} \qty{69.32}{EUR/t}. \textbf{7.} Gas, by \qty{0.33}{EUR/MWh} (89.35 against 89.68).
\textbf{8.} $(0.8525 - 0.3673) \times 10^6 = 485\,227$ tonnes. \textbf{9.} Above the \omterm{def:m3:emissions-and-environmental-markets:switch}{switching price},
\qty{69.32}{EUR/t}, ignoring each plant's other costs. \textbf{10.} It rises to \qty{125.53}{EUR/t}:
coal becomes the cheaper plant at \qty{70}{}.
\textbf{11.} With the TNAC at 1\,023\,494\,202, the reserve withholds 190\,494\,202 allowances from
auctions from September 2026 to August 2027: supply falls. \textbf{12.} Lower demand lowers the price
and raises the TNAC, which raises future intakes and tightens supply later. \textbf{13.} The marginal
plant's cost includes $0.8525 \times P_{\mathrm{CO_2}}$ per MWh, which the auction passes into the price.
\textbf{14.} Financial investors and funds, speculators, and intermediaries who sell forward to
compliance buyers; banks that hedge utilities. \textbf{15.} Nothing directly for its domestic power
sales, but imported electricity and energy-intensive goods now carry a carbon cost, which changes
competition. \textbf{16.} Its margin is a clean spread: the power price, fuel and carbon move together
in the \omterm{def:m3:power-markets-design:merit}{merit order}, and hedging one leg alone leaves the spread open. \textbf{17.} That the plant earns
little above fuel and carbon and will run less as the cap tightens: closure or conversion. \textbf{18.}
When gas is cheap, coal-to-gas switching is the cheapest large abatement, so the \omterm{def:m3:emissions-and-environmental-markets:switch}{switching price} acts as
a ceiling that the allowance price rarely exceeds while switching capacity remains. \textbf{19.}
\emph{Named result:} a \omterm{def:m3:emissions-and-environmental-markets:switch}{switching price} of \qty{69.32}{EUR/t}, and 6.584 million allowances (about
EUR~460.9 million at \qty{70}{EUR/t}) for the plan. \textbf{20.} The cost of the marginal tonne of
abatement needed to stay under the cap, now and in expectation (allowances can be banked).
\end{solution}

\begin{solution}{iq:m3:emissions-and-environmental-markets:1}
A tax fixes the price per tonne and lets emissions adjust; \omterm{def:m3:emissions-and-environmental-markets:cap}{cap-and-trade} fixes the quantity and lets the
price adjust through trade. Under uncertainty about abatement costs one fixes the outcome, the other the
cost.
\iqlookfor{price versus quantity instrument, and who bears uncertainty.}
\end{solution}

\begin{solution}{iq:m3:emissions-and-environmental-markets:2}
The power price minus the gas and carbon costs of one MWh from a gas plant. Negative when power is
cheaper than the plant's fuel plus allowances: in hours set by renewables or cheaper plants.
\iqlookfor{the formula with efficiency and emission factor.}
\end{solution}

\begin{solution}{iq:m3:emissions-and-environmental-markets:3}
It makes future auction supply decrease with the surplus: a high TNAC withdraws allowances, a low one
releases them. Supply becomes a function of the market's own holdings, which dampens persistent
surpluses and supports prices after demand shocks.
\iqlookfor{endogenous supply.}
\end{solution}

\begin{solution}{iq:m3:emissions-and-environmental-markets:4}
Dearer gas raises the \omterm{def:m3:emissions-and-environmental-markets:switch}{switching price}, so more coal runs and power emits more per MWh: demand for
allowances rises, and prices tend to rise; against that, high energy prices cut industrial output and
demand. The net depends on which dominates.
\iqlookfor{fuel switching as a demand channel, and the offsetting activity effect.}
\end{solution}

\begin{solution}{iq:m3:emissions-and-environmental-markets:5}
Quality risk (credits whose reduction did not happen or would have happened anyway), permanence risk
(forests that burn), reputational and legal risk (claims found misleading), and price risk if it must
replace invalidated credits.
\iqlookfor{additionality, permanence, claims.}
\end{solution}

\begin{solution}{iq:m3:emissions-and-environmental-markets:6}
A dispatch model: hourly demand, renewable output scenarios, the fleet with efficiencies and emission
factors, fuel and carbon price scenarios, and interconnection; run the \omterm{def:m3:power-markets-design:merit}{merit order} per hour and sum
emissions. Data: fleet registries, historical dispatch, weather. Calibrate on past verified emissions.
\iqlookfor{a merit-order dispatch model and scenarios, calibrated to verified emissions.}
\end{solution}
